% GENERATED FILE. Edit canonical lessons, then run npm run generate:books.
% canonical-source-hash: fnv1a64:df5d5726e52cf588
% canonical-lessons: RU-C70-monday, RU-C70-wednesday, RU-C70-thursday, RU-C70-friday, RU-C70-saturday

\chapter{Monday to Saturday}
\label{ch:ru-monday-to-saturday}
\hlchaptermodality{70}

\begin{chapteropening}
\textbf{By the end of this chapter:} \emph{I can name Monday, Wednesday, Thursday, Friday and Saturday.}

Complete the last lesson of chapter 70: I can name Monday, Wednesday, Thursday, Friday and Saturday.
\end{chapteropening}

\section[\texorpdfstring{\ru{понедельник} (ponedélnik)}{понедельник (ponedélnik)}]{\ru{понедельник} (ponedélnik) — Monday}

\label{lesson:RU-C70-monday}



\begin{quote}
Before the new one: say the Russian for a minute, then the Russian for time.
\end{quote}

\subsection*{You'll want to know: \ru{понедельник}}
\textbf{\ru{понедельник}} (\emph{ponedélnik}) — \textquotedblleft{}Monday\textquotedblright{}. \textquotedblleft{}After the week-day\textquotedblright{}: \textbf{\ru{неделя}} once meant Sunday.

\begin{cousinweb}
\textbf{\ru{понедельник}} is \emph{\ru{по}-}, \textquotedblleft{}after\textquotedblright{}, with \emph{\ru{неделя}}, which once meant \textquotedblleft{}Sunday\textquotedblright{}, and \emph{-\ru{ник}}: \textquotedblleft{}the day after Sunday\textquotedblright{} (Vasmer).
\end{cousinweb}

\begin{grammarlens}[title={What you've built}]
One more word for this chapter.
\end{grammarlens}

\subsection*{Your turn}
\begin{itemize}
\raggedright
  \item \emph{Say it:} \emph{ponedélnik}
  \item \emph{Say it:} \emph{ponedélnik}, once more
  \item \emph{Say it:} say \emph{ponedélnik}
  \item \emph{Recall:} say the Russian for a minute, then the Russian for time, then say \emph{ponedélnik} again
\end{itemize}

\begin{tcolorbox}[breakable,colback=teal!4,colframe=teal!35!black,title={Before you move on}]
What does \ru{понедельник} mean? (\textquotedblleft{}Monday\textquotedblright{}.) Say it once more.
\end{tcolorbox}
\section[\texorpdfstring{\ru{среда} (sredá)}{среда (sredá)}]{\ru{среда} (sredá) — Wednesday}

\label{lesson:RU-C70-wednesday}



\begin{quote}
Before the new one: say the Russian for time, then the Russian for Monday.
\end{quote}

\subsection*{You'll want to know: \ru{среда}}
\textbf{\ru{среда}} (\emph{sredá}) — \textquotedblleft{}Wednesday\textquotedblright{}. \textquotedblleft{}The middle\textquotedblright{} of the week.

\begin{cousinweb}
\textbf{\ru{среда}} — the Church Slavonic shape, says Vasmer, of native \emph{\ru{середа}}, \textquotedblleft{}middle\textquotedblright{}, the word in \emph{\ru{середина}}, \textquotedblleft{}the middle\textquotedblright{}: Wednesday is the middle of the week. It is the same root as \emph{\ru{сердце}}, \textquotedblleft{}heart\textquotedblright{}.
\end{cousinweb}

\begin{grammarlens}[title={What you've built}]
One more word for this chapter.
\end{grammarlens}

\subsection*{Your turn}
\begin{itemize}
\raggedright
  \item \emph{Say it:} \emph{sredá}
  \item \emph{Say it:} \emph{sredá}, once more
  \item \emph{Say it:} say \emph{sredá}
  \item \emph{Recall:} say the Russian for time, then the Russian for Monday, then say \emph{sredá} again
\end{itemize}

\begin{tcolorbox}[breakable,colback=teal!4,colframe=teal!35!black,title={Before you move on}]
What does \ru{среда} mean? (\textquotedblleft{}Wednesday\textquotedblright{}.) Say it once more.
\end{tcolorbox}
\section[\texorpdfstring{\ru{четверг} (chetvérg)}{четверг (chetvérg)}]{\ru{четверг} (chetvérg) — Thursday}

\label{lesson:RU-C70-thursday}



\begin{quote}
Before the new one: say the Russian for Monday, then the Russian for Wednesday.
\end{quote}

\subsection*{You'll want to know: \ru{четверг}}
\textbf{\ru{четверг}} (\emph{chetvérg}) — \textquotedblleft{}Thursday\textquotedblright{}. \textquotedblleft{}The fourth\textquotedblright{} day, from \textbf{\ru{четыре}}.

\begin{cousinweb}
\textbf{\ru{четверг}} $\leftarrow$ Proto-Slavic \emph{*četvĭrgŭ}, from the root of \emph{\ru{четыре}}, \textquotedblleft{}four\textquotedblright{} (Vasmer): the fourth day after Sunday.
\end{cousinweb}

\begin{grammarlens}[title={What you've built}]
One more word for this chapter.
\end{grammarlens}

\subsection*{Your turn}
\begin{itemize}
\raggedright
  \item \emph{Say it:} \emph{chetvérg}
  \item \emph{Say it:} \emph{chetvérg}, once more
  \item \emph{Say it:} say \emph{chetvérg}
  \item \emph{Recall:} say the Russian for Monday, then the Russian for Wednesday, then say \emph{chetvérg} again
\end{itemize}

\begin{tcolorbox}[breakable,colback=teal!4,colframe=teal!35!black,title={Before you move on}]
What does \ru{четверг} mean? (\textquotedblleft{}Thursday\textquotedblright{}.) Say it once more.
\end{tcolorbox}
\section[\texorpdfstring{\ru{пятница} (pyátnitsa)}{пятница (pyátnitsa)}]{\ru{пятница} (pyátnitsa) — Friday}

\label{lesson:RU-C70-friday}



\begin{quote}
Before the new one: say the Russian for Wednesday, then the Russian for Thursday.
\end{quote}

\subsection*{You'll want to know: \ru{пятница}}
\textbf{\ru{пятница}} (\emph{pyátnitsa}) — \textquotedblleft{}Friday\textquotedblright{}. \textquotedblleft{}The fifth\textquotedblright{} day, from \textbf{\ru{пять}}.

\begin{cousinweb}
\textbf{\ru{пятница}} is built on \emph{\ru{пять}}, \textquotedblleft{}five\textquotedblright{}: the fifth day of a week that begins on Monday (Vasmer); Old Russian said \emph{pjatŭkŭ}.
\end{cousinweb}

\begin{grammarlens}[title={What you've built}]
One more word for this chapter.
\end{grammarlens}

\subsection*{Your turn}
\begin{itemize}
\raggedright
  \item \emph{Say it:} \emph{pyátnitsa}
  \item \emph{Say it:} \emph{pyátnitsa}, once more
  \item \emph{Say it:} say \emph{pyátnitsa}
  \item \emph{Recall:} say the Russian for Wednesday, then the Russian for Thursday, then say \emph{pyátnitsa} again
\end{itemize}

\begin{tcolorbox}[breakable,colback=teal!4,colframe=teal!35!black,title={Before you move on}]
What does \ru{пятница} mean? (\textquotedblleft{}Friday\textquotedblright{}.) Say it once more.
\end{tcolorbox}
\section[\texorpdfstring{\ru{суббота} (subbóta)}{суббота (subbóta)}]{\ru{суббота} (subbóta) — Saturday}

\label{lesson:RU-C70-saturday}



\begin{quote}
Before the new one: say the Russian for Thursday, then the Russian for Friday.
\end{quote}

\subsection*{You'll want to know: \ru{суббота}}
\textbf{\ru{суббота}} (\emph{subbóta}) — \textquotedblleft{}Saturday\textquotedblright{}.

\begin{cousinweb}
\textbf{\ru{суббота}} — Vasmer supports a source in Byzantine Greek \emph{*sámbaton}; others prefer Balkan Latin \emph{*sambata}. Both come from Hebrew \emph{shabbat}, \textquotedblleft{}the day of rest\textquotedblright{}.
\end{cousinweb}

\begin{grammarlens}[title={What you've built}]
One more word for this chapter.
\end{grammarlens}

\subsection*{Your turn}
\begin{itemize}
\raggedright
  \item \emph{Say it:} \emph{subbóta}
  \item \emph{Say it:} \emph{subbóta}, once more
  \item \emph{Say it:} say \emph{subbóta}
  \item \emph{Recall:} say the Russian for Thursday, then the Russian for Friday, then say \emph{subbóta} again
\end{itemize}

\begin{tcolorbox}[breakable,colback=teal!4,colframe=teal!35!black,title={Before you move on}]
What does \ru{суббота} mean? (\textquotedblleft{}Saturday\textquotedblright{}.) Say it once more.
\end{tcolorbox}
